%==============================================================================
% Chapter 10: Conclusion
%==============================================================================
\chapter{Conclusion}
\label{chap:conclusion}

Modern enterprise cybersecurity is locked in an asymmetric conflict where adversaries deploy automated tooling while Security Operations Centers (SOCs) face manual triage bottlenecks, fragmented tooling, and severe analyst fatigue, exacerbated by commercial cloud AI copilots introducing grave data privacy risks, recurring token costs, and egress latency unsuitable for sovereign infrastructure. This project addressed these challenges by engineering and empirically validating an on-premise Sovereign Security Operations Center Assistant integrating Wazuh SIEM ingestion, Qdrant vector search, PostgreSQL persistence, and local Qwen3:8B inference via Ollama to establish a new operational standard for private, verifiable, and automated Tier-1 triage.

Our primary contributions encompass: first, a sovereign, air-gapped architecture executing entirely on local host hardware with zero cloud dependencies, eliminating data exfiltration risks and fulfilling GDPR and DPDP compliance; second, an entity-anchored retrieval breakthrough pairing deterministic observable extraction with Qdrant bitset payload filtering, elevating Recall@1 across 1,762 security chunks from 22.5\% to 90.0\% (+67.5 percentage points over standard dense RAG in 240.2~ms); third, an audited 120-run benchmark across 40 scenarios proving that while local 8B models exhibit strong parametric technique knowledge ($F_1 = 0.9028$), retrieval grounding is essential for vulnerability attribution (Log4Shell \texttt{CVE-2021-44228} achieved 1.00 $F_1$ vs.\ 0.0 for LLM-Only, 100\% non-CVE true-negative abstention, 75.0\% benign accuracy, and 97.3\% citation veracity); fourth, an exhaustive software verification suite validated by 264 automated tests (232 backend unit tests in 56.79~s and 32 frontend tests in 7.97~s, 100\% passing) and 6 regression tests enforcing alert namespacing and database isolation; and fifth, an interactive analyst workbench in React~19 featuring real-time SSE token streaming, interactive D3.js force-directed entity graphs, and automated remediation playbooks.

The practical operational implications of this research are substantial. Achieving a 95\% reduction in triage latency by completing structured entity extraction, knowledge retrieval, and incident synthesis in approximately 50 seconds transforms Tier-1 efficiency compared to historical 15--30 minute manual triage. Automating repetitive correlation of log lines, IP reputations, and MITRE tactics alleviates analyst cognitive burnout, freeing human defenders to focus on complex forensics and proactive threat hunting. Furthermore, demonstrating that fully grounded, citation-governed triage executes on commodity workstation GPUs (e.g., NVIDIA RTX 3060) democratizes advanced sovereign cyber defense for educational institutions and small-to-medium enterprises without cloud budgets.

Ultimately, artificial intelligence in cybersecurity is not an autonomous replacement for human judgment, but an indispensable force multiplier for defenders. By anchoring generative language models in authoritative threat knowledge bases and enforcing strict architectural boundaries of local data sovereignty and citation auditability, this research establishes that private, accurate, and explainable AI-assisted cyber defense is a practical, production-ready operational reality.
